% !TEX root = chapter-standalone.tex

\section{Feedback categories}
\label{sec:feedback-categories}



\begin{ctdefinition}[Feedback category]
    \label{def:feedback-cat}
    We call a \SY{strict monoidal category}~$\tup{\CatC, \mtimescat, \idmoncat, \braiding}$ a \emph{feedback category} if it is equipped with a \emph{feedback operator}: a family of functions
    \begin{equation}
        \Fb_{\Obja,\Objb}^\Objc\colon \HomSet{\CatC}{\Obja \mtimescatob \Objc}{\Objb\mtimescatob \Objc}\to \HomSet{\CatC}{\Obja}{\Objb},
    \end{equation}
    satisfying the following axioms:
    \begin{enumerate}
        \item \emph{Naturality in $\Obja$:} For any morphisms $\mora\colon \Obja\mtimescatob \Objc \mto \Objb\mtimescatob \Objc$ and $\morb \colon \Obja' \mto \Obja$,
              \begin{equation}\label{eq:feedback-naturality-1}
                  \Fb_{\Obja', \Objb}^\Objc ( (\morb \mtimescatmor \catidat\Objc) \mthen \mora) = \morb \mthen \Fb_{\Obja, \Objb}^\Objc (\mora)
              \end{equation}
        \item \emph{Naturality in $\Objb$:} For any morphisms $\mora\colon \Obja\mtimescatob \Objc \mto \Objb\mtimescatob \Objc$ and $\morb \colon \Objb \mto \Objb'$,
              \begin{equation}\label{eq:feedback-naturality-2}
                  \Fb_{\Obja, \Objb'}^\Objc ( \mora \mthen (\morb \mtimescatmor \catidat\Objc) ) = \Fb_{\Obja, \Objb}^\Objc (\mora) \mthen \morb
              \end{equation}
        \item \emph{Dinaturality in $\Objc$ with respect to isomorphisms:}
              For any morphism $\mora \colon \Obja  \mtimescatob \Objc \mto \Objb \mtimescatob \Objc'$ and any isomorphism $\morb \colon \Objc' \mto \Objc$,
              \begin{equation}\label{eq:feedback-dinaturality}
                  \Fb_{\Obja, \Objb}^{\Objc} (\mora \mthen (\catidat\Objb \mtimescatmor \morb)) = \Fb_{\Obja, \Objb}^{\Objc'}((\catidat\Obja \mtimescatmor \morb) \mthen \mora).
              \end{equation}

        \item \emph{Vanishing I:}
              For any morphisms~$\mora \colon \Obja \mtimescatob \idmoncat \mto \Objb \mtimescatob \idmoncat$ in \CatC,
              \begin{equation}
                  \label{eq:feedback-vanishing_1}
                  \Fb_{\Obja,\Objb}^\idmoncat (\mora)=  \mora.
              \end{equation}
        \item \emph{Vanishing II:}
              For any morphism~$\mora \colon (\Obja\mtimescatob \Objc) \mtimescatob \Objd \mto (\Objb\mtimescatob \Objc) \mtimescatob \Objd$ in \CatC,
              \begin{equation}
                  \label{eq:feedback-vanishing_2}
                  \Fb_{\Obja,\Objb}^\Objc\pars{
                      \Fb_{\Obja \mtimescatob \Objc, \Objb \mtimescatob \Objc}^\Objd(\mora)} = \Fb_{\Obja,\Objb}^{\Objc\mtimescatob \Objd}( \mora).
              \end{equation}
        \item \emph{Superposing:}
              For any morphism~$\mora\colon \Obja\mtimescatob \Objc \mto \Objb\mtimescatob \Objc$ in \CatC,
              \begin{equation}
                  \label{eq:feedback-superposing}
                  \Fb_{\Obje\mtimescatob \Obja,\Obje\mtimescatob \Objb}^{\Objc}( \catidat\Obje\mtimescatmor \mora  )=\catidat\Obje\mtimescatmor \Fb_{\Obja,\Objb}^\Objc(\mora).
              \end{equation}
    \end{enumerate}
\end{ctdefinition}

\begin{figure}[h!]
    \centering
    \begin{tikzpicture}[scale=0.7, every node/.style={transform shape},
        morbox/.style={draw=morphisms, fill=white, thick, minimum width=1.6cm, minimum height=2.2cm},
        wire/.style={objects, thick},
        wirelabel/.style={objects, font=\bfseries},
        tracebox/.style={draw=black, dashed, thick, fill=graphcolor, rounded corners=2pt},
    ]
        % Trace boundary (drawn first so it's behind everything)
        \fill[graphcolor, rounded corners=2pt] (-0.6,-2.4) rectangle (4.6,1.6);
        \draw[black, dashed, thick, rounded corners=2pt] (-0.6,-2.4) rectangle (4.6,1.6);

        % Morphism box f
        \node[morbox] (f) at (2,0) {};
        \node[morphisms, font=\large\itshape] at (2,0) {$f$};

        % X wire: input from left
        \draw[wire] (-1.8, 0.55) -- (f.west |- 0,0.55);
        \node[wirelabel, above] at (-1.2, 0.55) {$\Obja$};

        % Y wire: output to right
        \draw[wire] (f.east |- 0,0.55) -- (5.8, 0.55);
        \node[wirelabel, above] at (5.2, 0.55) {$\Objb$};

        % Z wire: input side (from feedback loop up to f)
        \draw[wire] (0.6, -1.8) -- (0.6, -0.55) -- (f.west |- 0,-0.55);
        \node[wirelabel, left] at (0.6, -1.1) {$\Objc$};

        % Z wire: output side (from f down to feedback loop)
        \draw[wire] (f.east |- 0,-0.55) -- (3.4, -0.55) -- (3.4, -1.8);
        \node[wirelabel, right] at (3.4, -1.1) {$\Objc$};

        % Feedback loop: bottom horizontal segment
        \draw[wire] (0.6, -1.8) -- (3.4, -1.8);
    \end{tikzpicture}
    \caption{The trace operator: given a morphism~$\mora \colon \Obja \mtimescatob \Objc \mto \Objb \mtimescatob \Objc$, its trace~$\Tr_{\Obja,\Objb}^\Objc(\mora) \colon \Obja \mto \Objb$ is obtained by ``closing the loop'' on the~$\Objc$ wire.}
    \label{fig:trace-operator}
\end{figure}

The diagrams in the figures \cref{fig:trace-naturality-x} to \cref{fig:trace-vanishing-ii} illustrate each of the axioms in the definition. Each axiom has an interpretation as a kind of compatibility condition. 

\begin{figure}[h!]
    \centering
    \includegraphics[width=10cm]{trace-naturality-x}
    \caption{Naturality in $\Obja$}
    \label{fig:trace-naturality-x}
\end{figure}

The axiom in \cref{fig:trace-naturality-x} is about compatibility of feedback with series composition via precomposition. 

\begin{figure}[h!]
    \centering
    \includegraphics[width=10cm]{trace-naturality-y}
    \caption{Naturality in $\Objb$}
    \label{fig:trace-naturality-y}
\end{figure}


The axiom in \cref{fig:trace-naturality-y} is about compatibility of feedback with series composition via postcomposition. 

\begin{figure}[h!]
    \centering
    \includegraphics[width=10cm]{trace-dinaturality-z}
    \caption{Dinaturality in $\Objc$}
    \label{fig:trace-dinaturality-z}
\end{figure}

The axiom in \cref{fig:trace-dinaturality-z} is about compatibility of feedback with series composition in the feedback variable; it is a kind of ``cyclicity'' condition. 

\begin{figure}[h!]
    \centering
    \includegraphics[width=6cm]{trace-vanishing-i}
    \caption{Vanishing I}
    \label{fig:trace-vanishing-i}
\end{figure}


The axiom in \cref{fig:trace-vanishing-i} is about compatibility of feedback with the monoidal unit. 

\begin{figure}[h!]
    \centering
    \includegraphics[width=10cm]{trace-vanishing-ii}
    \caption{Vanishing II}
    \label{fig:trace-vanishing-ii}
\end{figure}

The axiom in \cref{fig:trace-vanishing-ii} is about compatibility of feedback with the monoidal product. 



% Flush the pending float queue before the next mdframed environment.
% Without this, the seven consecutive [h!] figures above leave no space on the
% page and mdframed loops forever ("Not enough space on this page").
\FloatBarrier



\subsection{Basic examples}
\begin{ctdefinition}[Trace for $\RelL$ with concatenation as monoidal product]
    \label{def:tuple-REL-trace}
    Consider the symmetric strict monoidal category $\RelL$, where the monoidal product is defined via the concatenation operation $\cprod$ for lists of sets.
    Given a relation $\relA \colon \cObj{\setA, \setC}  \to \cObj{\setB, \setC}$, its \emph{trace} is defined as
    %
    \begin{equation}
        \label{eq:tupel-REL-tracedef}
        \begin{aligned}
            \Tr_{\cObj{\setA}, \cObj{\setB}}^{\cObj{\setC}}(\relA) = \{ & \tup{\ela, \elb} \in \cObj{\setA} \times \cObj{\setB} \mid \\
                                                                           & \exists \ \elc \in \cObj{\setC} : \tup{\ela \tupconcat \elc, \elb \tupconcat \elc} \in \relA \}
        \end{aligned}
    \end{equation}
\end{ctdefinition}


\subsection{Applied examples}

\todotext{J: I think the example below needs to be double-checked in comparison with the paper/notes by Roman on feedback categories -- there he quotients also by a relation that seems to give cyclicity of the feedback... but perhaps it amounts, in the end, to the same equivalence relation as the one we used (with math msc student's help...). UPDATE: I've checked: Roman's equivalence relation is equivalent to the one we used}



\begin{example}[Feedback for Mealy machines]
    \label{exa:mealy-feedback}
    The category $\MealyMod$ of Mealy machines up to isomorphism (\cref{def:MealyCat}) has a monoidal product.
    It is given, on objects, by the \SY{cartesian product} of sets, and, on morphisms, by running two machines side by side (compare \cref{exa:MealyModL-strict-monoidal}).
    A feedback operator for $\MealyMod$ is defined as follows.

    Let
    \begin{equation}
        \equivmora \colon \setA \cartprod \setC \mto \setB \cartprod \setC
    \end{equation}
    be a morphism, represented by a Mealy machine
    \begin{equation}
        \mora = \tupp{\setA \cartprod \setC,\ \prst,\ \setB \cartprod \setC,\ \prdyn_\mora,\ \prreadout_\mora},
    \end{equation}
    with set of states $\prst$, dynamics, and read-out
    \begin{equation}
        \prdyn_\mora \colon (\setA \cartprod \setC) \cartprod \prst \sto \prst,
        \qquad
        \prreadout_\mora \colon (\setA \cartprod \setC) \cartprod \prst \sto \setB \cartprod \setC.
    \end{equation}
    Since the set of outputs of $\mora$ is a \SY{cartesian product}, its read-out consists of two components,
    \begin{equation}
        \begin{aligned}
            \prreadout_\mora^{\setB} & \definedas \prreadout_\mora \mthen \proj_1 \colon (\setA \cartprod \setC) \cartprod \prst \sto \setB, \\
            \prreadout_\mora^{\setC} & \definedas \prreadout_\mora \mthen \proj_2 \colon (\setA \cartprod \setC) \cartprod \prst \sto \setC,
        \end{aligned}
    \end{equation}
    so that $\prreadout_\mora(\prinel, \prstel) = \tup{\prreadout_\mora^{\setB}(\prinel, \prstel), \prreadout_\mora^{\setC}(\prinel, \prstel)}$ for every input $\prinel \setin \setA \cartprod \setC$ and every state $\prstel \setin \prst$.

    We let
    \begin{equation}
        \Fb_{\setA, \setB}^{\setC}(\equivmora) \colon \setA \mto \setB
    \end{equation}
    be the morphism of $\MealyMod$ which is represented by the Mealy machine
    \begin{equation}
        \label{eq:mealy-feedback-machine}
        \tupp{\setA,\ \setC \cartprod \prst,\ \setB,\ \prdyn_{\Fb},\ \prreadout_{\Fb}}
    \end{equation}
    whose set of states is $\setC \cartprod \prst$, and whose dynamics and read-out are
    \begin{equation}
        \label{eq:mealy-feedback-dyn}
        \begin{aligned}
            \prdyn_{\Fb} \colon \setA \cartprod (\setC \cartprod \prst) & \sto \setC \cartprod \prst, \\
            \tup{\styleelements{a}, \tup{\styleelements{c}, \prstel}}                                        & \mapsto \tup{\prreadout_\mora^{\setC}(\tup{\styleelements{a}, \styleelements{c}}, \prstel),\ \prdyn_\mora(\tup{\styleelements{a}, \styleelements{c}}, \prstel)},
        \end{aligned}
    \end{equation}
    and
    \begin{equation}
        \label{eq:mealy-feedback-ro}
        \begin{aligned}
            \prreadout_{\Fb} \colon \setA \cartprod (\setC \cartprod \prst) & \sto \setB, \\
            \tup{\styleelements{a}, \tup{\styleelements{c}, \prstel}}                                            & \mapsto \prreadout_\mora^{\setB}(\tup{\styleelements{a}, \styleelements{c}}, \prstel).
        \end{aligned}
    \end{equation}
    In other words, the feedback operator $\Fb_{\setA, \setB}^{\setC}$ moves $\setC$ from the inputs and the outputs of the machine into its state.
    The $\setB$-component of the read-out of $\mora$ becomes the read-out of the new machine, and the $\setC$-component of the read-out of $\mora$ becomes part of the \emph{dynamics} of the new machine: it is no longer shown to the outside, but it is stored, and it is fed back to $\mora$ as the $\setC$-component of its input in the next step.

    Concretely, if the new machine is in the state $\tup{\styleelements{c}_{i-1}, \prsteln{i-1}}$ and receives the input $\styleelements{a}_i$, then it returns the output $\styleelements{b}_i$ and moves to the state $\tup{\styleelements{c}_i, \prsteln{i}}$, where
    \begin{equation}
        \label{eq:mealy-feedback-steps}
        \begin{aligned}
            \styleelements{b}_i          & = \prreadout_\mora^{\setB}(\tup{\styleelements{a}_i, \styleelements{c}_{i-1}}, \prsteln{i-1}), \\
            \styleelements{c}_i          & = \prreadout_\mora^{\setC}(\tup{\styleelements{a}_i, \styleelements{c}_{i-1}}, \prsteln{i-1}), \\
            \prsteln{i}  & = \prdyn_\mora(\tup{\styleelements{a}_i, \styleelements{c}_{i-1}}, \prsteln{i-1}).
        \end{aligned}
    \end{equation}
    So the value which $\mora$ emits on the wire $\setC$ in one step is the value which $\mora$ receives on the wire $\setC$ in the following step: this is feedback with a delay of one step.

    The morphism $\Fb_{\setA, \setB}^{\setC}(\equivmora)$ does not depend on the choice of the representative $\mora$.
    Indeed, if $\mapc \colon \prst \sto \prst'$ is an isomorphism of Mealy machines from $\mora$ to a machine $\mora'$, then it is compatible with both components of the read-outs, and therefore the function
    \begin{equation}
        \setC \cartprod \prst \sto \setC \cartprod \prst', \qquad \tup{\styleelements{c}, \prstel} \mapsto \tup{\styleelements{c}, \mapc(\prstel)},
    \end{equation}
    is an isomorphism of Mealy machines between the machines \cref{eq:mealy-feedback-machine} obtained from $\mora$ and from $\mora'$.

    If we combine the dynamics and the read-out of $\mora$ into a single function $(\setA \cartprod \setC) \cartprod \prst \sto (\setB \cartprod \setC) \cartprod \prst$, as in \cref{rem:mealy-single-function}, then the single function $\setA \cartprod (\setC \cartprod \prst) \sto \setB \cartprod (\setC \cartprod \prst)$ of the new machine is obtained from it simply by re-bracketing its arguments and its values.

    Finally, note that \cref{def:feedback-cat} asks for a \SY{strict monoidal category}, which $\MealyMod$ with the \SY{cartesian product} is not.
    The same construction applies to the strict version $\MealyModL$ of \cref{exa:MealyModL-strict-monoidal}, where it becomes even simpler: there, tuples are concatenated instead of nested, and no re-bracketing is needed.
\end{example}

\begin{remark}[Feedback and initial states]
    \label{rem:mealy-feedback-initial-states}
    The feedback operator of \cref{exa:mealy-feedback} makes use of the fact that our Mealy machines do not include an initial state.
    The set of states of the new machine is $\setC \cartprod \prst$, so an initial state for it would have to contain an element of $\setC$: an initial value on the wire that is fed back.
    The machine $\mora$ does not provide such a value, and there is no natural choice for it; the set $\setC$ may even be empty.
    For Mealy machines with initial state (\cref{def:MealyModst}), a feedback operation of this kind therefore needs an initial value on the wire $\setC$ as an additional piece of data.
\end{remark}
